\Needspace{12\baselineskip}
\section{References}

\begingroup
\setlength{\parskip}{0pt}
\titlespacing*{\subsection}{0pt}{1.0ex}{0.3ex}

\subsection*{Sacred Scripture}

{\fontsize{8.6}{9.6}\selectfont
\begin{itemize}[itemsep=0pt,parsep=0pt,topsep=0.05em]
\item Genesis 50:19--20; Job 2:10; Proverbs 16:9; Wisdom 11:21--26; Lamentations 3:37--38; Matthew 6:25--34 and 10:29--31; Acts 2:23; Romans 8:28; James 4:13--15; 1~Peter 5:6--7. Quoted from the Douay--Rheims (Challoner revision; public domain); chapter texts checked 2026-07-25 at \sourceurl{https://drbo.org/}{drbo.org}.
\end{itemize}}

\subsection*{Authoritative Catholic teaching and condemned propositions}

{\fontsize{8.6}{9.6}\selectfont
\begin{itemize}[itemsep=0pt,parsep=0pt,topsep=0.05em]
\item \emph{Catechism of the Catholic Church} 302--314 (divine providence; providence and secondary causes; providence and the scandal of evil) and 321--324 (in brief); official English text checked 2026-07-25 at \sourceurl{https://www.vatican.va/archive/ENG0015/__P19.HTM}{vatican.va}, Part One, Section Two, Chapter One, Article 1, Paragraph 4 (``The Creator'').
\item Innocent~XI, decree of the Holy Office of 28 August and constitution \emph{Coelestis Pastor} (1687), condemning sixty-eight propositions of Miguel de Molinos; Denzinger DS 2201--2269 (older numbering DB 1221--1288). Latin text with both numberings and the graded censure checked 2026-07-25 at \sourceurl{https://www.patristica.net/denzinger/}{patristica.net}; an English translation of the constitution's narrative sections read the same day at \sourceurl{https://www.papalencyclicals.net/innoc11/i11coel.htm}{papalencyclicals.net}. The Denzinger heading dates the constitution 20 November 1687; the 1913 \emph{Catholic Encyclopedia} dates it 2 November 1687. English renderings of the propositions in this article are working glosses; the Latin governs.
\item Innocent~XII, brief \emph{Cum alias ad apostolatus} (12 March 1699), condemning twenty-three propositions of François de Fénelon on the love of God; Denzinger DS 2351--2374 (DB 1327--1349), with the censure at DS 2374. Latin checked 2026-07-25 at \sourceurl{https://www.patristica.net/denzinger/}{patristica.net}. The censure lists the propositions as rash, scandalous, ill-sounding, offensive to pious ears, pernicious in practice, and erroneous, respectively; it does not include a charge of heresy.
\end{itemize}}

\subsection*{Thomistic framework}

{\fontsize{8.6}{9.6}\selectfont
\begin{itemize}[itemsep=0pt,parsep=0pt,topsep=0.05em]
\item Thomas Aquinas, \emph{Summa theologiae} I, q.~19, aa.~6 and 9 (the antecedent and consequent will; whether God wills evils); q.~22, aa.~1--4 (providence as the plan of things ordered to an end; its universal extension; immediacy of design and mediated execution; contingency preserved); q.~103, aa.~5--8 (the universality of divine government; government through intermediaries; nothing outside or resisting the order of government); q.~105, aa.~4--5 (God moves the created will without coercing it; God works in every agent without displacing it); q.~116, aa.~1--4 (fate); II-II, q.~21, a.~1 (presumption); II-II, q.~136, aa.~1--5 (patience). English from the public-domain translation of the Fathers of the English Dominican Province, checked 2026-07-25 at \sourceurl{https://www.newadvent.org/summa/}{newadvent.org}; Latin checked the same day at \sourceurl{https://www.corpusthomisticum.org/iopera.html}{corpusthomisticum.org}. Where the witnesses differ, the Latin governs.
\end{itemize}}

\subsection*{The classic text and its sources}

{\fontsize{8.6}{9.6}\selectfont
\begin{itemize}[itemsep=0pt,parsep=0pt,topsep=0.05em]
\item \emph{Trustful Surrender to Divine Providence: The Secret of Peace and Happiness}, by Fr.~Jean Baptiste Saint-Jure, S.J., and St.~Claude de la Colombière, S.J., trans.\ Paul Garvin (Charlotte, N.C.: TAN Books, 2012 printing; nihil obstat John P. Sullivan, imprimatur Francis Cardinal Spellman, 8 November 1961; first published 1961 by Alba House as \emph{The Secret of Peace and Happiness}; reissued by St.~Raphael Editions 1978, 1980, and 1983 with TAN; copyright 1980 St.~Raphael Editions). In copyright. Front matter, publisher's note, table of contents, and opening pages checked 2026-07-25 in the publisher's own preview file at \sourceurl{https://tanbooks.com/content/Preview_TrustfulSurrender.pdf}{tanbooks.com}.
\item Jean-Baptiste Saint-Jure, S.J., \emph{De la connaissance et de l'amour du fils de Dieu notre seigneur Jésus Christ}, vol.~2 (Paris: Périsse frères, 1877), Book~III, ch.~VIII, §§ I--X (pp.~206--282). Public domain; scan and OCR checked 2026-07-25 at the \sourceurl{https://archive.org/details/delaconnaissance02sain}{Internet Archive}.
\item \emph{Œuvres du R.~P. Claude de la Colombière, de la Compagnie de Jésus}, vol.~5 (1832): sermons \emph{Sur la soumission à la volonté de Dieu} (p.~65), \emph{Sur la confiance en Dieu} (p.~85), \emph{Sur la prière} (p.~102), \emph{Sur les adversités} (p.~218). Public domain; checked 2026-07-25 at the \sourceurl{https://archive.org/details/uvresdurpclauded05laco}{Internet Archive}. Volumes 1--4 and the 1684 \emph{Réflexions chrétiennes} were also searched.
\end{itemize}}

\subsection*{The tradition of abandonment}

{\fontsize{8.6}{9.6}\selectfont
\begin{itemize}[itemsep=0pt,parsep=0pt,topsep=0.05em]
\item Francis de Sales, \emph{Treatise on the Love of God} (1616), Books~VIII (chs.~I--III) and IX (chs.~I, IV--VI, VIII, XIV), in the public-domain English translation of Henry Benedict Mackey, O.S.B.; text checked 2026-07-25 at the \sourceurl{https://ccel.org/ccel/desales/love.all.html}{Christian Classics Ethereal Library}.
\item Attributed to Jean-Pierre de Caussade, S.J. (see the historical section above), \emph{Abandonment; or, Absolute Surrender to Divine Providence}, revised and corrected by Henri Ramière, S.J., trans.\ Ella McMahon from the eighth French edition (New York, Cincinnati and Chicago: Benziger Brothers, 1887; imprimatur Michael Augustine, Archbishop of New York, 15 February 1887): Ramière's preface, pp.~3--20, and Book~I, ch.~I, pp.~31--33. Public domain; checked 2026-07-25 at the \sourceurl{https://archive.org/details/abandonmentorabs00caus}{Internet Archive}.
\item Same work, \emph{Abandonment to Divine Providence}, third English edition, trans.\ E.~J. Strickland from the tenth complete French edition (Exeter: Catholic Records Press, 1921), introduction by Dom Arnold, O.S.B. Public domain; consulted 2026-07-25 at the \sourceurl{https://archive.org/details/divineprovidence00causuoft}{Internet Archive} for edition identity only.
\item Dominique Salin, S.J., ``The Treatise on Abandonment to Divine Providence,'' \emph{The Way} 46/2 (April 2007): 21--36. In copyright; quoted briefly. Checked 2026-07-25 at the journal's posted PDF, \sourceurl{https://www.theway.org.uk/back/462Salin.pdf}{theway.org.uk}. Reports the manuscript history, Ramière's editorial interventions, and the scholarship of Michel Olphe-Galliard, S.J., and Jacques Gagey (\emph{L'Abandon à la providence divine d'une dame de Lorraine au XVIIIe siècle}, Grenoble: Jérôme Million, 2001), neither of which was consulted directly.
\end{itemize}}

\subsection*{Historical background}

{\fontsize{8.6}{9.6}\selectfont
\begin{itemize}[itemsep=0pt,parsep=0pt,topsep=0.05em]
\item ``Miguel de Molinos,'' \emph{The Catholic Encyclopedia}, vol.~10 (New York: Robert Appleton Company, 1911); public domain; checked 2026-07-25 at \sourceurl{https://www.newadvent.org/cathen/10441a.htm}{newadvent.org}. Used as a reported witness for Molinos's biography, the publication and approbations of \emph{La Guía espiritual}, the 1685 arrest, and the September 1687 sentence.
\item The beatification (1929) and canonisation (31 May 1992) of Claude de la Colombière are reported at the level of standard reference works, checked 2026-07-25; no act of the Holy See was consulted directly, and no argument depends on the dates.
\end{itemize}}

\endgroup
